\documentclass[11pt,a4paper]{article}
\usepackage[T1]{fontenc}
\usepackage{lmodern,microtype}
\usepackage[margin=20mm,headheight=14pt]{geometry}
\usepackage{amsmath,amssymb,mathtools,booktabs,tabularx,array}
\usepackage{enumitem,xcolor,tikz,pgfplots,caption,listings}
\usetikzlibrary{arrows.meta,positioning,calc,matrix,fit,decorations.pathreplacing}
\pgfplotsset{compat=1.18}
\usepackage{fancyhdr,natbib,bibentry}
\definecolor{courseblue}{HTML}{1F4E8C}
\definecolor{courseteal}{HTML}{247C78}
\definecolor{muted}{HTML}{586575}
\definecolor{softblue}{HTML}{EDF3FA}
\definecolor{softteal}{HTML}{EAF5F1}
\definecolor{courseorange}{HTML}{A85A29}
\usepackage[colorlinks=true,linkcolor=courseblue,urlcolor=courseblue]{hyperref}
\hypersetup{pdftitle={From Text to Recurrent Models: Lectures 01--04},pdfauthor={CS40008.01, Fudan University}}
\pagestyle{fancy}\fancyhf{}
\lhead{\small CS40008.01 / Fall 2026}
\rhead{\small Lecture notes / first draft}
\cfoot{\small\thepage}
\renewcommand{\headrulewidth}{0.3pt}
\setlength{\parindent}{0pt}\setlength{\parskip}{5pt}
\setlist[itemize]{leftmargin=1.3em,itemsep=2pt,topsep=2pt}
\setlist[enumerate]{leftmargin=1.5em,itemsep=2pt,topsep=2pt}
\setcounter{secnumdepth}{1}
\captionsetup{font=small,labelfont=bf,skip=5pt,hypcap=false}
\renewcommand{\thefigure}{F\ifnum\value{figure}<10\relax0\fi\arabic{figure}}
\newcolumntype{Y}{>{\raggedright\arraybackslash}X}
\newcommand{\R}{\mathbb{R}}
\newcommand{\BOS}{\textsc{bos}}
\newcommand{\EOS}{\textsc{eos}}
\newcommand{\UNK}{\textsc{unk}}
\newcommand{\softmax}{\operatorname{softmax}}
\newcommand{\sigmoid}{\sigma}
\newcommand{\readingref}[1]{\hyperlink{reading:#1}{#1}}
\newcommand{\doi}[1]{\href{https://doi.org/#1}{DOI}}
\newcommand{\paperentry}[7]{%
  \par\addvspace{7pt}\hypertarget{reading:#1}{}%
  {\color{courseblue}\bfseries #1\quad #2}\hfill{\footnotesize\bfseries #3}\par
  {\footnotesize\color{muted}#4\par}{\small\bibentry{#5}\par}%
  \textbf{Why it matters.} #6\par\textbf{Read for.} #7\par}
\newcommand{\nextpage}[1]{\clearpage\subsection*{#1}}
\newcommand{\takeaway}[1]{\par\textcolor{courseteal}{\textbf{Takeaway.}} #1\par}
\newcommand{\selfcheck}[1]{\par\textcolor{courseblue}{\textbf{Check your understanding.}} #1\par}
\newcommand{\fig}[2]{\begin{center}\input{figures/#1.tex}\captionof{figure}{#2}\label{fig:#1}\end{center}}
\tikzset{box/.style={draw=courseblue,fill=softblue,rounded corners=2pt,align=center,minimum height=8mm,font=\small},
  tbox/.style={box,draw=courseteal,fill=softteal},
  arr/.style={-{Latex[length=2mm]},thick,courseblue},
  flow/.style={-{Latex[length=2mm]},thick,courseteal},
  lab/.style={font=\small,align=center}}
\pgfplotsset{every axis/.append style={font=\small,tick label style={font=\footnotesize},
  grid=major,grid style={gray!15},axis line style={gray!60},legend style={draw=none,font=\footnotesize}}}
\lstset{language=Python,basicstyle=\ttfamily\small,keywordstyle=\color{courseblue},
  commentstyle=\color{muted},showstringspaces=false,columns=fullflexible,
  frame=single,rulecolor=\color{gray!30},breaklines=true,aboveskip=8pt,belowskip=8pt}
\begin{document}
\bibliographystyle{plainnat}\nobibliography{references}
\thispagestyle{empty}
{\small\color{courseblue}\bfseries CS40008.01 / FUDAN UNIVERSITY / FALL 2026\par}
\vspace{12pt}
{\fontsize{26}{31}\selectfont\bfseries From Text to Recurrent Models\par}
\vspace{6pt}
{\large Lecture notes for Lectures 01--04\par}
{\color{muted}First draft / October 1, 2026\par}

Language modeling connects several ideas that initially look separate:
representing text, estimating probabilities, learning vectors, and remembering
context. These notes develop that connection from a sequence of small,
inspectable examples. The common task is to predict the next token from the
information available so far.

\fig{course-map}{The four lectures share a prediction task. Tokenization determines
the symbols; the model determines how context becomes a probability distribution.}

\section*{A route through the notes}
\begin{tabularx}{\textwidth}{@{}r Y r@{}}\toprule
& Topic & Page\\\midrule
1 & \hyperref[sec:task]{Language modeling as a common task} & \pageref{sec:task}\\
2 & \hyperref[sec:tokens]{From text to tokens} & \pageref{sec:tokens}\\
3 & \hyperref[sec:counts]{Count-based language models} & \pageref{sec:counts}\\
4 & \hyperref[sec:evaluation]{Evaluating and sampling language models} & \pageref{sec:evaluation}\\
5 & \hyperref[sec:embeddings]{Learning token representations} & \pageref{sec:embeddings}\\
6 & \hyperref[sec:neural]{From token IDs to a trainable neural LM} & \pageref{sec:neural}\\
7 & \hyperref[sec:rnn]{Recurrent language models} & \pageref{sec:rnn}\\
8 & \hyperref[sec:lstm]{LSTMs and the limits of recurrent memory} & \pageref{sec:lstm}\\
& \hyperref[sec:review]{Review, notation, and sources} & \pageref{sec:review}\\
& \hyperref[sec:readings]{Guide to 15 papers before the Transformer} & \pageref{sec:readings}\\\bottomrule
\end{tabularx}

\textbf{How to read.} Follow each worked example with pencil and paper.
The checks ask you to explain a boundary condition, identify an assumption,
or recompute a value. The course notebooks provide executable companions.
All diagrams are editable vector figures; measured examples are labeled with
their source and limitations.

\textbf{Scope.} The revised Lecture 04 ending develops RNNs and LSTMs.
Attention begins in Lecture 05. The final section prepares that transition
by identifying the limitations that remain after gated memory.

\clearpage
\section{Language modeling as a common task}\label{sec:task}
\subsection*{From text applications to a probability model}
A translation, a summary, and a continuation all contain sequences of text.
A language model assigns probabilities to token sequences. This makes it
possible to score candidate sequences and generate new ones. A high model
probability expresses what the model expects; factual accuracy and success
on a task still need their own checks.

Let $x_1,\ldots,x_T$ be tokens from a finite output vocabulary
$\mathcal V$ of size $K$. We first fix the tokenizer and the treatment of
document boundaries. The chain rule gives the exact identity
\[
 p_\theta(x_{1:T})=\prod_{t=1}^{T}p_\theta(x_t\mid x_{<t}).
\]
This identity makes no short-context assumption. An $n$-gram, a feedforward
network, and an RNN differ in how they approximate each conditional.
Every conditional must be nonnegative and sum to one over $\mathcal V$.

\subsection*{The learning objective}
Write $P$ for the data distribution and $Q_\theta$ for the model. On a
fixed sequence sample space, with $Q_\theta$ positive wherever $P$ is positive,
\[
 D_{\mathrm{KL}}(P\Vert Q_\theta)
 =\mathbb E_P[\log P(X)]-\mathbb E_P[\log Q_\theta(X)].
\]
The first term is independent of $\theta$, so minimizing this divergence
amounts to maximizing expected log likelihood. We cannot sum over the
unknown data distribution directly. Instead, training uses observed documents
and minimizes their empirical negative log likelihood:
\[
 \mathcal L(\theta)=-\frac{1}{N_{\mathrm{pred}}}
 \sum_{d}\sum_{t\in\mathcal S_d}\log p_\theta(x_t^{(d)}\mid x_{<t}^{(d)}).
\]
Here $\mathcal S_d$ specifies the scored positions, including an end token
if the evaluation convention scores it. $N_{\mathrm{pred}}$ is their total
count. During training this denominator is constant with respect to the
parameters. Averaging each document separately and then averaging documents
would give shorter documents more weight per token.

\subsection*{Worked example: probability is a product, loss is a sum}
Suppose the model assigns probabilities $1/2$, $1/4$, and $1/2$ to the
three scored events in a short document. The document probability is $1/16$,
its negative log probability is $\log 16$, and its mean token loss is
$\log(16)/3\approx0.9242$ nats. Taking logs converts a long product of
small numbers into a sum and makes the contribution of each prediction visible.

\textbf{Boundaries matter.} A beginning marker supplies context but need
not be an output symbol. An end marker can be a predicted event. Including
the probability of termination is part of defining a distribution over
variable-length documents; a generation length cap is a separate practical rule.

\takeaway{Fix the representation and scoring events before comparing models.
Then ask how each model turns the allowed context into a normalized distribution.}
\selfcheck{If one target receives probability zero, what happens to the document
probability and its loss? The answers are zero and $+\infty$. Explain why
adding an end token changes the scored-event count.}
\textbf{Reading connection.} \readingref{R01} connects uncertainty with
coding; \readingref{R04} replaces discrete context tables with shared learned features.

\clearpage
\section{From text to tokens}\label{sec:tokens}
\subsection*{Text processing makes modeling decisions}
A tokenizer maps text to integer IDs; its decoder maps those IDs back to
text. The complete pipeline includes earlier decisions about decoding,
normalization, and splitting. A later decoder cannot reconstruct distinctions
that an earlier step has discarded. Split documents into training,
development, and test sets before fitting vocabulary or preprocessing choices.

\fig{text-pipeline}{Every arrow has a contract. Normalization may intentionally
lose information; lossless tokenization applies to the text after that step.}

A visible symbol, a Unicode code point, and a UTF-8 byte are different
units. The precomposed letter \texttt{\'e} is one code point and two bytes.
The visually similar sequence \texttt{e} followed by a combining acute accent
uses two code points and three bytes. Rendering alone cannot identify the
stored sequence.

\fig{unicode-units}{Two encodings of an accented letter. Hexadecimal byte values
describe storage; grapheme appearance describes what a reader sees.}

NFC normalization can compose canonically equivalent sequences. NFKC also
applies compatibility mappings, which may remove distinctions useful to a
task. Lowercasing is another transformation: it maps both \texttt{A} and
\texttt{a} to the same string. None is a universal requirement for modeling.
Document the intended transformation and evaluate on text processed the same way.

\fig{boundaries}{Extraction can omit punctuation and spaces. A lossless partition
covers the complete string and can also constrain which adjacent pieces may merge.}

For the input \texttt{hi, 3!}, extracting alphabetic runs produces only
\texttt{hi}. A partition into \texttt{hi}, \texttt{,}, a space, \texttt{3},
and \texttt{!} preserves every character. A regular expression can implement
either behavior; the word ``regex'' alone says nothing about reversibility.

\selfcheck{If a tokenizer round-trips lowercased text, does it also recover the
original capitalization? No: test the complete preprocessing pipeline separately.}

\nextpage{Learning a vocabulary with byte-pair encoding}
Word occurrences count positions, while word types count distinct forms.
For \texttt{low low lower}, whitespace splitting gives three occurrences
and two types; UTF-8 storage uses 13 bytes, including two spaces. Token IDs
are arbitrary labels, not numerical measurements of meaning.

Byte-pair encoding (BPE) begins from a base alphabet and repeatedly adds a
token formed by merging a frequent adjacent pair. In the course's byte-level
example the base vocabulary contains all 256 byte values. Each training
sequence has a frequency weight; boundaries prevent merges between sequences.
The toy corpus is \texttt{low}$\times5$ and \texttt{lower}$\times2$.

\fig{bpe-merges}{Two weighted merges. Both initial winning pairs have count seven;
the course's deterministic tie rule selects \texttt{l o} first.}

The initial weighted length is $5(3)+2(5)=25$. The pair
$(\texttt{l},\texttt{o})$ occurs once in each of seven weighted words, so
its merge reduces the total to 18. Merging $(\texttt{lo},\texttt{w})$
reduces it to 11. The new IDs are 256 for \texttt{lo} and 257 for
\texttt{low}; two merges give 258 vocabulary entries.

\textbf{Training algorithm.}
\begin{enumerate}
\item Count adjacent pairs inside each current sequence and multiply by
the sequence's corpus frequency.
\item Select the most frequent pair, resolving ties deterministically.
Assign its concatenation a new token ID and a merge rank.
\item Replace non-overlapping occurrences of that pair, then recount.
Stop after the requested number of merges or when no eligible pair remains.
\end{enumerate}
Counting and replacement are different operations. In \texttt{aaa}, the
pair \texttt{a a} occurs at two overlapping positions, but one replacement
can consume only two of the three symbols. For \texttt{aaa}$\times2$ and
\texttt{bc}$\times2$, its weighted pair count is four while the token total
decreases only from ten to eight.

\takeaway{Vocabulary growth is learned from training text. Deterministic
ties, allowed boundaries, and overlap handling are part of the algorithm.}
\selfcheck{Recompute the second merge counts after the first replacement.
Why would merging across spaces change what the vocabulary can represent?}
\textbf{Reading connection.} \readingref{R03} applies BPE to subword units
within words. Its character-based setup differs from this byte-level example.

\nextpage{Encoding new text and choosing a vocabulary}
Encoding freezes the learned merge list. Begin with the new text's base
symbols, find an eligible adjacent pair with the earliest learned rank,
merge it, and repeat. Do not count new-text frequencies to learn new rules:
the model's embedding rows already correspond to a fixed vocabulary.

\fig{frozen-bpe}{The learned \texttt{lo} and \texttt{low} merges apply to
\texttt{lowest}; the remaining bytes have no learned merges in this example.}

The resulting IDs are $[257,101,115,116]$. Decoding concatenates each
token's byte string and then decodes the complete UTF-8 byte stream.
A single token can end inside a multibyte character, so decoding tokens
independently may fail. Byte coverage ensures that valid input bytes can be
represented; it does not ensure efficient segmentation for every language.
Special tokens such as \BOS{} and \EOS{} have separate control meanings
and need an explicit policy when similar strings appear in ordinary text.

\fig{vocabulary-tradeoff}{Left: larger vocabularies increase table rows while often
shortening sequences. Right: held-out Chinese token counts from the L01 toy
experiment, using the same two sentences at every merge budget.}

At 0, 8, and 32 merges, mixed English/Chinese training gives 54, 48, and 35
tokens on those Chinese sentences. English-only training gives 54 each time.
The tokenizer still covers the text, but its merges offer no compression there.
These small constructed corpora demonstrate a mechanism; they do not rank
tokenizers on broad language-model quality.

For a table with $K$ rows and width $d$, storage grows as $Kd$. A materialized
output distribution at $BT$ positions has $BTK$ entries. Shorter sequences
may reduce position-dependent work, while a larger vocabulary increases
table and readout costs. The resulting tradeoff depends on the model and
text distribution; token compression alone cannot settle it.

\takeaway{Fit the tokenizer on training data, freeze it for evaluation, and
report both representation behavior and downstream modeling results.}
\selfcheck{Why can a tokenizer with fewer tokens have higher per-token perplexity?
Section 4 will separate the prediction unit from the amount of text scored.}

\clearpage
\section{Count-based language models}\label{sec:counts}
\subsection*{Approximate the context, then estimate probabilities}
The chain rule permits the full prefix. An $n$-gram approximation keeps
only the last $n-1$ tokens:
\[
 p(x_t\mid x_{<t})\approx p(x_t\mid x_{t-n+1:t-1}).
\]
A unigram has no token context, a bigram has one context token, and a
trigram has two. ``Trigram'' counts the predicted token as well as its context.
The approximation concerns conditional independence, not the factorization itself.

\fig{chain-context}{A prediction for the final token can use increasingly long
histories. The full-prefix chain rule remains exact; the short histories impose a model assumption.}

For a history $h$, let $C(h,w)$ count training occurrences followed by $w$,
and let $C(h)=\sum_{v\in\mathcal V}C(h,v)$. Maximum likelihood gives
\[
 \widehat p(w\mid h)=\frac{C(h,w)}{C(h)},\qquad C(h)>0.
\]
To see why, minimize $-\sum_w C(h,w)\log p_w$ subject to $\sum_wp_w=1$.
The Lagrange multiplier condition is $-C(h,w)/p_w+\lambda=0$ for observed
events; normalization gives $\lambda=C(h)$.

\fig{count-probability}{The course example has three occurrences of history
\texttt{I}: two followed by \texttt{am}, one by \texttt{do}. Counts and probabilities have different units.}

We extract adjacent pairs separately inside each sentence, including the
chosen boundary markers. Joining sentences first could invent a pair from
the last word of one sentence to the first word of the next. In the
L02/L03 example, \BOS{} is context-only and \EOS{} is a possible target.
The row for \texttt{I} assigns $2/3$ to \texttt{am} and $1/3$ to
\texttt{do}; all other targets receive zero under unsmoothed MLE.

\selfcheck{Which count belongs in the denominator: occurrences of the predicted
word or occurrences of the history? Explain why each conditional row must
sum to one, even if a plotted table puts histories in columns instead.}

\nextpage{Sparse counts, smoothing, and unseen histories}
A zero training count is often weak evidence that an event is impossible.
Yet unsmoothed MLE assigns zero probability and infinite loss to such an
event. Additive smoothing reserves probability for every symbol in a
fixed output vocabulary:
\[
 p_\delta(w\mid h)=\frac{C(h,w)+\delta}{C(h)+\delta K},\qquad\delta>0.
\]
The denominator is the sum of the numerators, so normalization is explicit.
For an unseen history, $C(h)=0$ gives the uniform distribution $1/K$.
This is a defined fallback, unlike the undefined MLE ratio $0/0$.

\fig{smoothing}{An original three-symbol example with counts $(2,1,0)$.
Add-one smoothing changes the row to $(1/2,1/3,1/6)$; interpolation uses a
chosen uniform fallback with weight $0.3$.}

For a mixture of normalized component models,
\[
 p(w\mid h)=\sum_{r=0}^{n}\lambda_r p_r(w\mid h),\qquad
 \lambda_r\ge0,\quad\sum_r\lambda_r=1,
\]
the result is normalized because summation over $w$ commutes with the finite
mixture. Components can use different history lengths; $p_0$ can be a
uniform floor. Each component must define its behavior for unseen histories.
Tune mixture weights on development text, then hold them fixed for the test set.

\textbf{A useful distinction: frequency versus continuation diversity.}
Imagine a word that appears 100 times, always after the same preceding
word, and another that appears only ten times after ten different words.
The latter has appeared in more kinds of contexts. Kneser--Ney smoothing
uses continuation statistics for its lower-order distribution, rather than
simply recycling unigram frequency. \readingref{R02} develops this idea;
the recap requires the distinction, not a full implementation.

\textbf{Unknown words and output support.} A word tokenizer may map unseen
forms to \UNK{}. A byte tokenizer can represent their bytes, but the language
model can still encounter unseen byte or subword contexts. Input-only
markers should not silently enter the output denominator. Vocabulary size
means the number of outcomes the model actually predicts in this example.

\takeaway{A language model needs a complete normalized rule for every
history encountered at evaluation or generation time.}
\selfcheck{For $\delta=1$ and counts $(2,1,0)$, verify the unseen symbol gets
$1/6$. What happens as $\delta\to\infty$? The row approaches the uniform distribution.}

\clearpage
\section{Evaluating and sampling language models}\label{sec:evaluation}
\subsection*{One total loss, several units}
Let $S=-\sum\log p_\theta(x_t\mid x_{<t})$ be total negative log likelihood
in nats, $N$ the number of scored predictions, and $B_{\mathrm{utf8}}$ the
number of evaluated text bytes. The same score can be reported as
\[
 \ell=\frac SN,\quad \mathrm{BPT}=\frac{S}{N\log2},\quad
 \mathrm{PPL}=e^{S/N},\quad\mathrm{BPB}=\frac{S}{B_{\mathrm{utf8}}\log2}.
\]
Perplexity is the inverse geometric mean of the probabilities assigned to
observed targets. For a uniform distribution on ten digits, every target
has probability $1/10$, so every nonempty sequence has perplexity ten.
For a general model it is an effective branching factor, not an actual
count of equally likely options at every position.

\fig{metric-units}{Start from a common numerator. The denominator determines
whether a quantity is per token or per byte; exponentiation produces perplexity.}

\fig{tokenizer-metrics}{A constructed comparison with four total bits of loss
on the same four-byte text. Fewer scored tokens increase per-token perplexity
here; bits per byte stays one.}

The figure holds total loss fixed to isolate units. Real tokenizers also
change the model's events and fitted probabilities, so BPB can change too.
Cross-tokenizer BPB comparisons require identical evaluated text,
normalization, document boundaries, and special-token scoring conventions.

\textbf{End tokens affect the numerator and token count.} If a one-byte
document \texttt{a} has probability $1/2$ for its character and $1/2$ for
\EOS{}, then $N=2$, $S=2\log2$, and $B_{\mathrm{utf8}}=1$. Perplexity
is two and BPB is two. The boundary token contributes loss without adding
a literal byte to the evaluated text.

Use training loss to diagnose fitting, development loss to make choices,
and a fixed test set for final comparisons. Task-based evaluation asks a
different question: a model can improve average token loss while becoming
worse on a particular task or subgroup.

\selfcheck{Why does averaging the perplexities of two batches generally differ
from exponentiating their combined mean loss? Accumulate total loss and
prediction counts first, especially when batches have different sizes.}

\nextpage{Sampling, data filters, and experimental conclusions}
To sample a categorical distribution, draw $u$ uniformly from $[0,1)$ and
find the interval containing it in the cumulative probabilities. For a
language model, sample a token, append it to the context, recompute the
distribution, and repeat until \EOS{} or a stated length cap.

\fig{sampling}{A categorical draw for $(1/2,1/3,1/6)$. A draw at $u=0.7$
selects the second interval; generation repeats the same operation with updated context.}

For an interpolated model, either sample from the combined distribution or
first draw a component according to its mixture weight and then sample
from that component. Both implement the same probability law. Sampling
from raw counts while reporting the likelihood of a smoothed mixture
would describe two different models. Temperature and top-$k$ truncation
also change the generation distribution and should be labeled.

\fig{pipeline-experiments}{Two model-based data workflows in L02. A filter makes
a selection from scored documents; a sample/retrain loop changes the data
used by the next model. Neither diagram establishes a benefit on its own.}

\textbf{A filter score is a signal.} A reference LM may give lower loss to
text resembling its own training distribution. Selecting low-loss documents
can therefore change topic, language, and style. The course's reference
trigram is a teaching analogue of perplexity-based filtering; a threshold
does not certify truth, usefulness, or broad data quality.

\textbf{Read the budget with the curve.} In the saved L02 sample/retrain
demonstration, held-out BPB rises from 1.299 to 1.940 over five rounds.
The original training corpus has 6,076,505 tokens; the final synthetic
round has 100,681. Each round samples 2,000 documents capped at 128 tokens,
so length caps and changing token budgets are confounded with replacing
real data by samples. The observation is deterioration in this run, not
an isolated causal estimate of one mechanism.

An informative follow-up would hold total training tokens, vocabulary,
evaluation text, and training procedure fixed, then compare real-only,
synthetic-only, and mixed data across several seeds. This is an experiment
design, not a result reported by the notes.

\takeaway{Scoring, sampling, and selecting training data are distinct
operations. State the distribution and the comparison each one uses.}
\selfcheck{If a sampler stops at a length cap without \EOS{}, should the
result be recorded as naturally terminated? Keep the two stopping reasons distinct.}

\clearpage
\section{Learning token representations}\label{sec:embeddings}
\subsection*{From context counts to association}
Words that occur in similar contexts can acquire related representations.
The choice of ``context'' matters: neighboring words, a document, and a
syntactic relation define different statistics. In the lecture's constructed
word--context table, \texttt{tea} and \texttt{coffee} share \texttt{drink}
and \texttt{hot}, while \texttt{car} occurs with \texttt{drive}.

Let $X_{wc}$ be the count for word $w$ and context $c$, and
$Z=\sum_{w,c}X_{wc}$. Pointwise mutual information compares the observed
joint frequency with the frequency expected under independence:
\[
 \operatorname{PMI}(w,c)=\log_2\frac{X_{wc}Z}{X_{w\cdot}X_{\cdot c}},\qquad
 \operatorname{PPMI}(w,c)=\max(0,\operatorname{PMI}(w,c)).
\]
Here the row and column totals are positive. An unobserved pair with positive
marginals has PMI $-\infty$ and is assigned PPMI zero.

\fig{counts-ppmi}{The L03 invented count table and its PPMI transform. Counts
sum to 35. The transformed values measure association in bits, not probability.}

For \texttt{tea}/\texttt{drink}, the joint count is eight, the row total is
14, and the column total is 15. Therefore
\[
 \operatorname{PMI}(\texttt{tea},\texttt{drink})
 =\log_2\frac{8\cdot35}{14\cdot15}
 =\log_2(4/3)\approx0.4150\text{ bits}.
\]
The zero entries after PPMI conflate unobserved and nonpositively associated
pairs. This choice is part of the representation, not a lossless rewrite of counts.

\subsection*{Compression and alternative objectives}
For a matrix $M$, truncated SVD forms $M_r=U_r\Sigma_rV_r^\top$.
It minimizes squared reconstruction error among matrices of rank at most
$r$. A word table can use $U_r\Sigma_r$, producing a dense vector for each
row; the scaling convention should be stated because it affects geometry.

GloVe instead fits weighted log counts using vectors and biases:
\[
 J=\sum_{X_{wc}>0}f(X_{wc})
 \bigl(u_w^\top v_c+b_w+\widetilde b_c-\log X_{wc}\bigr)^2.
\]
The weight function reduces the influence of selected count ranges.
This objective differs from applying SVD to PPMI. Both produce vectors,
but what they preserve is determined by their input statistics and loss.

\selfcheck{Why can a frequent word have low PMI with a frequent context?
The independence baseline is already large. See \readingref{R07} for GloVe's objective.}

\nextpage{Predictive embeddings and one gradient by hand}
CBOW predicts a center word from its surrounding context; skip-gram predicts
surrounding contexts from a center word. These are representation-learning
tasks. A symmetric context window may contain future words, so it should
not be mistaken for the prefix-only conditioning of a causal language model.

\fig{cbow-skipgram}{Prediction direction changes the training examples.
Word2vec uses distinct input and output roles even when both index the same vocabulary.}

Negative sampling replaces a vocabulary-wide prediction objective with
observed-versus-sampled pair discrimination. For one positive context $c_+$
and one sampled negative context $c_-$, write
\[
 \ell=-\log\sigma(w^\top c_+)-\log\sigma(-w^\top c_-),\qquad
 \sigma(a)=\frac1{1+e^{-a}}.
\]
The sigmoid values do not sum to one across vocabulary items. This loss
is therefore not a normalized next-token negative log likelihood.

\fig{negative-gradient}{The positive score is one and the negative score is
$0.5$. Gradient descent increases the positive score when the center vector is held fixed.}

Using $w=(1,0.5)$, $c_+=(0.5,1)$, and $c_-=(1,-1)$ gives
$\sigma(1)\approx0.7311$, $\sigma(0.5)\approx0.6225$, and
$\ell\approx1.2873$. Differentiating the positive term gives
\[
 \nabla_{c_+}\ell=(\sigma(w^\top c_+)-1)w
 \approx(-0.2689,-0.1345).
\]
At step size $\eta=0.5$, the positive context becomes approximately
$(0.6345,1.0672)$. With $w$ fixed, its score rises from 1 to 1.1681.
The negative-context gradient is $\sigma(w^\top c_-)w$; descent reduces
that score. The center vector receives contributions from both pairs.

\takeaway{Specify the prediction task and objective before interpreting
the vectors. The sampling distribution also influences what the model learns.}
\selfcheck{At zero scores, each binary term contributes $\log2$. What is the
initial loss with one positive and five negatives? It is $6\log2\approx4.1589$.}
\textbf{Reading connection.} \readingref{R05} introduces CBOW/skip-gram;
\readingref{R06}, Section 2.2, gives the negative-sampling objective.

\nextpage{Static vectors, contextual states, and subword composition}
An embedding table is a learned collection of rows. Looking up the ID for
\texttt{bank} returns the same row every time that ID occurs. Contextual
computation can transform that row together with other inputs into a
different state for ``river bank'' and ``bank loan.'' Learning a good table
does not by itself supply that context computation.

\fig{representation-types}{A token row, a contextual state, and a retrieval
vector have different inputs and uses. A retrieval representation may summarize
an entire passage and be trained for matching.}

Cosine similarity compares the directions of two nonzero vectors:
\[
 \operatorname{cos}(u,v)=\frac{u^\top v}{\|u\|_2\|v\|_2}.
\]
Nearest neighbors and vector analogies are useful inspections of geometry.
They depend on the corpus, objective, normalization, and chosen examples.
A plausible neighbor list does not establish language-model likelihood,
factual accuracy, or performance on a downstream task.

\fig{subword-composition}{An illustrative fastText-style composition sums
overlapping character $n$-gram vectors. BPE segmentation supplies a sequence
of token IDs. The two operations use subwords in different ways.}

In a character $n$-gram representation, a word can be represented by a
word-specific vector plus vectors for its subword features. Overlapping
features can share statistical strength between related forms, and known
features can help represent a previously unseen word. The exact boundary
symbols, $n$-gram lengths, hashing, and normalization belong to the chosen
implementation. The figure only illustrates the composition principle.

In a BPE-based language model, the tokenizer instead emits separate subword
tokens. Each token has a lookup row; the model then processes their ordered
sequence. Character-feature composition and BPE can address related coverage
problems, but they define different model inputs.

\textbf{What carries forward.} Modern neural LMs still begin with lookup
tables and learn output scores. The additional question is how the model
forms a state from its available context. We next make the lookup, target
alignment, objective, and parameter updates explicit.

\selfcheck{If two occurrences have the same token ID, when must their vectors
be identical? At a deterministic table lookup. After a context-dependent
network, their states can differ.}

\clearpage
\section{From token IDs to a trainable neural LM}\label{sec:neural}
\subsection*{Lookup and next-token alignment}
Let $E\in\R^{K_{\mathrm{in}}\times d}$ be the input table. For token ID $i$,
lookup returns $E[i]$. Equivalently, for a one-hot row vector $q_i$,
$q_iE=E[i]$. The one-hot construction explains the operation; an embedding
lookup avoids materializing the mostly zero tensor. IDs have integer dtype,
while the trainable weights and selected vectors are floating-point tensors.

For a batch of $B$ sequences with $T$ input positions, IDs have shape
$(B,T)$ and embeddings have shape $(B,T,d)$. The L03 example uses a table
with ten rows and width four. The table shape and the input shape are
different objects; neither is inferred from the numerical magnitude of an ID.

\fig{token-loss}{The neural bigram computation keeps one input token per
prediction. Lookup changes the representation; it does not extend the context.}

For the sequence $[1,2,3,4,5]$, use inputs $[1,2,3,4]$ and targets
$[2,3,4,5]$. A second row $[5,6,7,8,9]$ gives inputs $[5,6,7,8]$ and
targets $[6,7,8,9]$. There are eight predictions, with last targets 5 and 9.
For a vocabulary of size ten, logits have shape $(2,4,10)$.

An untied neural bigram computes
\[
 e_t=E[x_t],\qquad z_t=W_oe_t+b_o,\qquad
 p(x_{t+1}\mid x_t)=\softmax(z_t),
\]
where $W_o\in\R^{K\times d}$. In the general case, $K_{\mathrm{in}}$ and
$K$ need not be equal: an input-only boundary symbol may occupy a lookup
row without being an allowed output. L03's ten-symbol toy uses matching sizes.

Flatten logits to $(BT,K)$ and targets to $(BT)$ in the same order before
computing cross-entropy. Sliced tensors may be non-contiguous; a reshape
can preserve their logical order while making a copy if needed. Mixing batch
and time orders would silently train on the wrong labels.

\takeaway{Write down shape, dtype, device, and meaning at every step.
Correct dimensions alone do not prove correct target alignment.}
\selfcheck{For $B=2$, $T=4$, $d=4$, and $K=10$, how many logits exist?
There are 80; only eight are the logits of the observed target classes.}

\nextpage{Stable probabilities and a connected computation graph}
For one target $y$ and logits $z\in\R^K$, the loss is
\[
 \ell(z,y)=-z_y+\log\sum_{j=1}^{K}e^{z_j}.
\]
Let $a=\max_jz_j$. Subtracting $a$ inside the exponentials gives the stable
log-sum-exp identity
\[
 \ell(z,y)=-(z_y-a)+\log\sum_j e^{z_j-a}.
\]
All exponent arguments are nonpositive, avoiding overflow from a large
common offset. For $z=(1000,1001)$ and the second class as target,
$\ell=\log(1+e^{-1})\approx0.3133$ rather than an expression involving
overflowing $e^{1000}$ and $e^{1001}$.

\fig{loss-graph}{Autograd follows the operations connecting parameters to the
loss. Detaching the loss breaks that connection even when its printed value is unchanged.}

The gradient with respect to logits has a simple form:
\[
 \frac{\partial\ell}{\partial z_j}=p_j-\mathbf1[j=y].
\]
It increases pressure on the target score and decreases pressure on the
others. Summing these derivatives gives zero because probabilities sum to
one. Adding a common constant to all logits leaves probabilities unchanged.

With all logits zero, $p_j=1/K$ and the loss is exactly $\log K$.
For $K=10$, this is about 2.3026 nats. Random logits generally do not yield
that exact value; use an intentional zero-logit setup for the check.

\begin{lstlisting}
# logits: (B, T, K); targets: (B, T), integer IDs
loss = F.cross_entropy(
    logits.reshape(-1, logits.shape[-1]),
    targets.reshape(-1),
)
loss.backward()
\end{lstlisting}
Pass raw logits to this loss, which performs the stable normalization
internally. Calling \texttt{softmax} first changes the function being
optimized. Calling \texttt{loss.item()} returns a Python number for
logging; it cannot carry the computation graph. Calling \texttt{detach()}
similarly removes the path needed for this backward call.

\selfcheck{If training loss is numerically finite but every parameter gradient
is absent, what would you inspect first? Trace the computation graph and
the gradient-recording context before changing the learning rate.}

\nextpage{A feedforward model can use more than one token}
Consider ``red key'' and ``blue key.'' If the next word should be
\texttt{opens} in the first context and \texttt{closes} in the second,
a bigram model cannot separate them: its input is \texttt{key} in both cases.
An embedding with more dimensions still receives the same input ID.

\fig{window-model}{An ordered two-token context distinguishes the two prefixes.
The hidden layer mixes the concatenated embeddings before producing vocabulary logits.}

A neural probabilistic language model uses a fixed window of $m$ tokens.
With column-vector conventions,
\[
 u=\operatorname{concat}(E[x_{t-m}],\ldots,E[x_{t-1}]),\quad
 h=\tanh(W_hu+b_h),\quad z=W_oh+b_o.
\]
Here $u\in\R^{md}$, $W_h\in\R^{H\times md}$, and
$W_o\in\R^{K\times H}$. Concatenation preserves slot order. The nonlinear
activation allows behavior beyond a single affine map of the concatenated
inputs; composing affine layers without an activation would collapse to
one affine map of $u$.

The L04 corpus contains two separate sequences:
\begin{center}
\texttt{BOS red key opens EOS}\qquad\texttt{BOS blue key closes EOS}.
\end{center}
For $m=2$, each supplies three windows. For the first sentence they are
$(\BOS,\texttt{red})\to\texttt{key}$,
$(\texttt{red},\texttt{key})\to\texttt{opens}$, and
$(\texttt{key},\texttt{opens})\to\EOS$. The second sentence supplies
the analogous blue/closes examples. No window crosses a document boundary.

\begin{tabularx}{\textwidth}{@{}l Y@{}}\toprule
Stage & Shape for a two-example probe, $m=2,d=4,H=8,K=7$\\\midrule
Context IDs & $(2,2)$\\
Lookup / concatenate & $(2,2,4)$ / $(2,8)$\\
Hidden / logits & $(2,8)$ / $(2,7)$\\\bottomrule
\end{tabularx}

This toy allocates all seven listed symbols as both input and output rows,
including \BOS{}, although \BOS{} never appears as a target. That convention
explains its $K=7$ baseline. A model with \BOS{} excluded from the output
would need a different output mapping and denominator.

\takeaway{Shared learned representations improve generalization across
contexts, but a fixed window still has a hard history limit. See \readingref{R04}.}
\selfcheck{Which parameter matrix grows when the context window doubles while
$d,H,K$ stay fixed? The input-to-hidden matrix $W_h$ doubles its column count.}

\nextpage{Training, tiny-batch fitting, and evaluation}
Gradient descent adjusts parameters to reduce a chosen objective. A
training step must compute a connected loss, differentiate it, and apply
an update. Gradients normally accumulate, so clear old gradients when
each step is intended to use only the current batch.

\begin{lstlisting}
model.train()
for step in range(200):
    optimizer.zero_grad(set_to_none=True)
    logits = model(train_contexts)
    loss = F.cross_entropy(logits, train_targets)
    loss.backward()
    optimizer.step()
\end{lstlisting}
This loop assumes the model, optimizer, and sentence-local training windows
from the course notebook have already been constructed. The call order is
the transferable idea; the batch and objective determine what it learns.

\fig{training-curves}{Saved L03 and L04 teaching runs, redrawn from separate
assets. Left: neural bigram, $K=10,d=4$, eight pairs, seed 0. Right: two-token
NPLM, $K=7,d=4,H=8$, six windows, seed 7. Both use 200 full-batch SGD updates
with learning rate $0.5$ on CPU. These curves are not a controlled model comparison.}

L03's mean loss reaches about 0.0187 nats, while L04's reaches about
0.0070. Their vocabularies, batches, architectures, and seeds differ.
Both examples establish that these implementations can fit their small
compatible training sets. Neither curve measures generalization to new documents.

An incompatible tiny batch can expose a model limitation. If a bigram sees
the same token with two equally frequent targets, its best conditional
distribution splits probability between them. Its optimum for those two
events has loss $\log2$, not zero. A failed fit is therefore not automatically
an optimizer bug: first check whether the model's inputs distinguish the labels.

For evaluation, switch behavior with \texttt{model.eval()} and separately
disable gradient recording with \texttt{torch.no\_grad()}. The first affects
modules such as dropout; the second controls graph construction. Neither
performs an optimizer step. The L03 unseen-ID probe tests coverage of selected
input rows; it is a diagnostic, not a representative held-out corpus.

\selfcheck{A training loss falls while held-out loss rises. Which observation
supports fitting, and which warns about generalization? Report both with their splits.}

\nextpage{Weight tying and a memory ledger}
If the same $K$ symbols occupy input and output rows and the state width is
$d$, the output matrix can share the input table: $W_o=E$. Then
$z=Eh+b_o$. This saves a table, but changes the gradient paths. Compatible
row meanings and widths are essential; equal tensor shapes alone are insufficient.

\fig{weight-tying}{Untied tables have independent parameters. Tying makes one
parameter receive the sum of its input-lookup and output-readout gradients.}

For an untied lookup, a row unused as input has no gradient through that
lookup. With tying, it may receive a gradient through output normalization
even if its token never appeared in the input batch. In a finite-logit
softmax, every output class generally contributes. Do not count the same
shared parameter twice when reporting model size.

\fig{memory-ledger}{Payload for one $32{,}000\times512$ table. Each fp32 tensor
of this shape uses 62.5 MiB; two Adam moment tensors account for 125 MiB.
Activations and temporary allocations require additional memory.}

The table has $P=32{,}000\cdot512=16{,}384{,}000$ parameters.
At two bytes per parameter its bf16 payload is 31.25 MiB. If parameters,
gradients, and two Adam moment tensors all use fp32, their combined payload is
$4P+4P+8P=16P$ bytes, or 250 MiB. This is an explicit storage model;
mixed-precision implementations can store different copies and dtypes.

Output logits add a different cost. Materializing a $(B,T,K)$ tensor needs
$BTK$ elements. A dense output projection from width $d$ costs approximately
$2BTdK$ forward FLOPs when one multiply-add counts as two. In the L03 toy,
$B=2,T=4,d=4,K=10$ gives 80 logits and about 640 projection FLOPs.
This excludes softmax, backward computation, and the optimizer update.

Allocated embedding rows, tokenizer entry counts, and maximum token ID are
not always identical. Use the actual matrix shape for storage, and check
that every emitted ID falls within its valid row range.

\takeaway{Separate unique parameters, tensor payload, intermediate tensors,
and peak device memory. A parameter count alone does not specify a training budget.}
\selfcheck{If weights are tied, do output logits disappear? No. Sharing removes
one parameter table, while the per-position prediction computation remains.}

\clearpage
\section{Recurrent language models}\label{sec:rnn}
\subsection*{Carry a state instead of fixing a window}
A fixed-window neural LM cannot distinguish prefixes whose final $m$
tokens agree. To let earlier tokens influence a prediction, a recurrent
network carries a state forward. At each step it combines the current
embedding with the previous state:
\[
 e_t=E[x_t],\qquad
 h_t=\tanh(W_xe_t+W_hh_{t-1}+b_h),\qquad
 z_t=W_oh_t+b_o.
\]
The state has width $H$, so $W_x\in\R^{H\times d}$,
$W_h\in\R^{H\times H}$, and $W_o\in\R^{K\times H}$.
The output $\softmax(z_t)$ predicts $x_{t+1}$. The same parameters are
used at every position, while the state changes with the sequence.

\fig{recurrent-state}{Unrolling a recurrent network exposes the sequence of
computations. Repeated boxes share parameters; the hidden states are different tensors.}

There is no fixed architectural cutoff at $m$ previous tokens in this
recurrence. An earlier input can affect later states through repeated
updates. This describes a possible dependency, not a guarantee that useful
information will be retained or learned. A finite-width, finite-precision
state must preserve relevant distinctions through those updates.

\subsection*{Worked scalar recurrence}
To isolate the state update, choose $h_0=0$, $W_x=1$, $W_h=0.5$,
$b_h=0$, and scalar inputs $(e_1,e_2,e_3)=(1,0,0)$. Then
\[
 h_1=\tanh(1)\approx0.7616,\quad
 h_2=\tanh(0.5h_1)\approx0.3634,\quad
 h_3=\tanh(0.5h_2)\approx0.1797.
\]
The first input still affects the third state, although its effect has
shrunk. These are hand-chosen parameters illustrating a mechanism, not
the output of a trained language model.

At a new independent document, reset or explicitly initialize the state.
Carrying a previous document's state would change the conditioning context.
For padded batches, exclude padding positions from the loss unless they
are intentionally part of the modeled event space. Reset policy and loss
masking are data conventions, not details to leave implicit.

\takeaway{Recurrence shares a transition rule across time. It trades a
fixed input window for a repeatedly updated summary of the prefix.}
\selfcheck{Does making a batch larger remove the dependence of $h_t$ on
$h_{t-1}$ within one sequence? No; batch parallelism and time dependence are separate.}
\textbf{Reading connection.} \readingref{R08} develops learned temporal
representations; \readingref{R11} applies recurrent states to language modeling.

\nextpage{Backpropagation through time}
Training an RNN differentiates the unrolled computation graph. A parameter
can affect the loss directly at the current step and indirectly through
many later states. Backpropagation through time (BPTT) applies the ordinary
chain rule to all these paths, then sums contributions to each shared parameter.

\fig{bptt-paths}{Forward states connect adjacent time steps. Backward signals
from a later loss pass through earlier states. A truncation boundary blocks
the backward connection while a numerical state can still be carried forward.}

For the vanilla recurrence, hold input embeddings fixed and define the
local recurrent Jacobian
\[
 J_t=\frac{\partial h_t}{\partial h_{t-1}}
 =\operatorname{diag}(1-h_t\odot h_t)W_h.
\]
The influence of an earlier state on a later one is a product:
\[
 \frac{\partial h_t}{\partial h_s}=J_tJ_{t-1}\cdots J_{s+1},\qquad s<t.
\]
The order matters because matrices do not generally commute. A loss at
step $t$ sends its derivative through this product to states before $t$.
When differentiating a shared weight, sum over the steps where that weight
was used as well as over the losses affected by each use.

\textbf{A small graph to trace.} With three recurrent steps and losses at
all three outputs, the first state receives contributions from all three
losses. The third state receives a contribution only from the third loss
in this graph. The transition weights are shared, so their gradient
accumulates contributions at steps one, two, and three.

\subsection*{Truncation changes the learning path}
For long streams, truncated BPTT divides computation into chunks. The
final state of one chunk may initialize the next, while its graph is
detached at the boundary. The next chunk sees earlier information in the
forward state, but its loss cannot change parameters through the detached
earlier computation. The forward context and backward credit-assignment
horizon are therefore different quantities.

For a first-order RNN, ordinary full BPTT stores information proportional
to sequence length for the backward pass. Checkpointing can trade extra
computation for less storage, but it does not alter the mathematical
gradient when the recomputation is exact. Truncation instead omits paths
and generally changes that gradient.

\selfcheck{If a clue is carried across a detached boundary, can the current
loss teach the earlier transition how to store that clue through that path?
No. It can still update parameters through their uses in the current chunk.}

\nextpage{Why long dependencies are difficult to learn}
Repeated multiplication can make sensitivity very small or very large.
For compatible matrix norms,
\[
 \left\|\frac{\partial h_t}{\partial h_s}\right\|
 \le\prod_{k=s+1}^{t}\|J_k\|.
\]
If every factor is bounded by a constant below one, the bound decreases
exponentially with distance. Factors larger than one permit amplification,
but do not by themselves prove explosion in every direction. Singular
directions, nonlinearities, and the actual state trajectory matter.

\fig{recurrent-gradients}{Exact scalar powers illustrate repeated multiplication.
The curves are not measured RNN gradients: real gradients involve matrix
products and contributions from several paths.}

At 50 steps, $0.9^{50}\approx0.00515$ and $1.1^{50}\approx117.39$.
The scalar example makes the issue visible without assuming that an RNN's
Jacobian is constant. A saturated $\tanh$ also has a small derivative,
so a model can represent a dependency in principle while receiving a weak
learning signal for it. \readingref{R09} studies the difficulty of learning
long dependencies with gradient descent.

\subsection*{What gradient clipping does}
For a gradient vector $g$ and threshold $\tau>0$, global norm clipping uses
\[
 \widetilde g=
 \begin{cases}
 g,&\|g\|_2\le\tau,\\
 \tau g/\|g\|_2,&\|g\|_2>\tau.
 \end{cases}
\]
For $g=(3,4)$ and $\tau=2$, the clipped vector is $(1.2,1.6)$.
It preserves direction and caps the norm of the gradient supplied to the
optimizer. With plain SGD it also bounds that update norm by $\eta\tau$;
adaptive optimizers can rescale coordinates afterward. Clipping does not
restore a learning signal that has already vanished.

\takeaway{Long-range learning depends on both forward memory and backward
sensitivity. Neither a long theoretical context nor gradient clipping
guarantees successful credit assignment.}
\selfcheck{Would setting every recurrent weight to one guarantee stable gradients?
No: the relevant factors are full local Jacobians, including nonlinear derivatives.}
\textbf{Supporting reading.} \href{https://proceedings.mlr.press/v28/pascanu13.html}{Pascanu,
Mikolov, and Bengio (2013)} analyzes these training difficulties and clipping.

\clearpage
\section{LSTMs and the limits of recurrent memory}\label{sec:lstm}
\subsection*{Separate stored content from exposed state}
A vanilla RNN replaces its state through a nonlinear transformation at
every step. LSTM introduces an additional cell state and learned gates
that control retaining, writing, and exposing information. The cell state
$c_t$ and hidden state $h_t$ each have width $H$ in the formulation below.

Let $q_t=[e_t;h_{t-1}]\in\R^{d+H}$. The modern cell used in these notes is
\begin{align*}
 f_t&=\sigma(W_fq_t+b_f), & i_t&=\sigma(W_iq_t+b_i),\\
 o_t&=\sigma(W_oq_t+b_o), & g_t&=\tanh(W_gq_t+b_g),\\
 c_t&=f_t\odot c_{t-1}+i_t\odot g_t,
 &h_t&=o_t\odot\tanh(c_t).
\end{align*}
Here the four gate/candidate matrices have shape $H\times(d+H)$.
Within this cell, $W_o,b_o$ name the output gate; a separate vocabulary
readout acts on $h_t$ to predict the next token.

\fig{lstm-cell}{A cell with a forget gate and no peephole connections. The top
path carries the cell state; multiplication nodes apply gates elementwise.}

\begin{itemize}
\item The \textbf{forget gate} $f_t$ controls how much previous cell content survives.
\item The \textbf{input gate} $i_t$ controls how much candidate content $g_t$ is written.
\item The \textbf{output gate} $o_t$ controls how much transformed cell content is exposed as $h_t$.
\end{itemize}
Sigmoid gates lie between zero and one, independently for each coordinate.
They are not a categorical distribution: $f_t+i_t$ need not equal one.
The candidate can be negative. The cell state accumulates content through
an additive update and is not itself restricted to $[-1,1]$.

\textbf{Historical convention.} \readingref{R10}, the 1997 LSTM paper,
introduced the memory-cell idea before this forget-gate formulation.
The extension is documented by \href{https://sferics.idsia.ch/pub/juergen/FgGates-NC.pdf}{Gers,
Schmidhuber, and Cummins (2000)}. Distinguishing the two avoids attributing
all modern equations to the original cell.

\selfcheck{Why keep $c_t$ as well as $h_t$? The model can retain cell content
while controlling how much of it reaches the prediction and next gate computations.}

\nextpage{A scalar LSTM update and its gradient path}
Choose $c_{t-1}=0.8$, $f_t=0.9$, $i_t=0.2$, $g_t=0.5$, and $o_t=0.75$.
These hand-chosen values make each operation visible:
\[
 c_t=0.9(0.8)+0.2(0.5)=0.82,\qquad
 h_t=0.75\tanh(0.82)\approx0.5063.
\]
The previous memory contributes $0.72$ and the candidate contributes $0.10$.
The hidden state exposes a bounded transform of their sum.

\fig{gate-arithmetic}{Idealized gate settings illustrate retain, overwrite,
and hide operations. Exact zero/one gates are limiting cases of sigmoid gates.}

If the gates and candidate are held fixed, increasing $c_{t-1}$ by $0.1$
increases $c_t$ by $0.09$. Along this direct cell-state path,
\[
 \left.\frac{\partial c_t}{\partial c_{t-1}}\right|_{f_t,i_t,g_t\ \mathrm{fixed}}
 =\operatorname{diag}(f_t).
\]
Across several steps, the contribution through only the direct cell path
is an elementwise product of forget gates. Gates near one can preserve
this contribution better than repeated saturated nonlinear transformations.

\subsection*{Why this is not the full gradient}
The gates depend on $h_{t-1}$, which depends on earlier cells and gates.
There are therefore additional paths from earlier states to later losses.
A product of forget gates isolates one path; it does not equal the full
recurrent derivative in general. The output gate and $\tanh(c_t)$ also
affect the sensitivity of a prediction to the stored content.

Even on the direct path, a forget value of 0.9 repeated 50 times yields
about 0.00515. Gates closer to one help retention, but successful learning
also requires that the right content be written, preserved, and exposed
at the right time. An LSTM improves the architecture's options for learning
memory; it does not eliminate all gradient or optimization difficulties.

\textbf{A useful alternative.} A GRU combines state and gating differently,
using reset and update gates. \readingref{R12} provides the historical
encoder--decoder context for that design. For this recap, understand the
shared purpose of controlling recurrent memory before comparing full equations.

\takeaway{Treat the additive memory path as a useful mechanism with stated
assumptions. Separate its local derivative from the total gradient of the network.}
\selfcheck{Can $h_t$ be near zero while $c_t$ retains a nonzero value? Yes:
a nearly closed output gate can hide stored content from the exposed state.}

\nextpage{What recurrent memory still cannot make easy}
\subsection*{Sequential computation remains}
The standard RNN or LSTM update needs the previous state. Even when all
training tokens are known, state $h_t$ cannot be evaluated by the usual
recurrence until $h_{t-1}$ is available. Matrix operations within a step
and different sequences in a batch can run in parallel; the within-sequence
chain remains a source of limited parallelism.

\fig{sequential-computation}{Two sequences can be processed alongside each
other, but each row has a left-to-right state dependency. Knowing the next
input does not supply the next hidden state.}

\subsection*{A state must retain the distinctions needed later}
A fixed-width state can be sufficient for some tasks, such as tracking a
small set of conditions. It need not preserve every detail of the prefix.
However, when useful information is spread across a long sequence, later
predictions depend on what survived earlier updates. Long-range credit
assignment must also teach the network which details to retain.

\fig{encoder-bottleneck}{A classic encoder--decoder interface passes only a
final encoder representation to the decoder. This particular design requires
that representation to carry the source information needed for every output.
Decoder token inputs are omitted.}

In the sequence-to-sequence model of \readingref{R13}, an LSTM encoder
produces a fixed-dimensional representation for an LSTM decoder. Unlike
an ordinary next-token LM on one stream, the task maps one sequence to
another. The interface makes the burden concrete: details needed at
different output positions must all remain available through that final
representation. This bottleneck is a property of the interface, not a
claim that every recurrent architecture exposes only one source vector.

The remaining concerns are related but distinct: sequential computation,
learning across long gradient paths, and retaining useful distinctions in
a finite state. Gates address part of the memory problem; they do not
remove the serial dependency or automatically change the encoder--decoder interface.

\takeaway{The next question is how a prediction could consult earlier
information more directly. Lecture 05 begins with learned alignment and
attention, using \readingref{R14} and \readingref{R15}.}
\selfcheck{Which limitation would remain if a recurrent model remembered
every needed detail perfectly? Its standard state updates would still depend on earlier steps.}

\clearpage
\section*{Review, notation, and sources}\label{sec:review}
\subsection*{Explain the complete modeling path}
Given a new text sample, you should be able to identify what preprocessing
preserves, how tokens are formed, which context each prediction uses, how
its probability is normalized, and how a learning signal reaches parameters.
The following questions connect the sections.

\begin{tabularx}{\textwidth}{@{}l Y@{}}\toprule
Topic & A question to answer without running code\\\midrule
Tokenization & Why can byte coverage coexist with poor compression for a language?\\
Counts & How do unseen histories differ from unseen continuations of a known history?\\
Evaluation & What changes in perplexity when the token unit changes?\\
Embeddings & What distinguishes pair discrimination from normalized next-token likelihood?\\
Neural LMs & Why cannot a wider token table alone distinguish two prefixes ending in the same token?\\
Training & What does successful tiny-batch fitting establish, and what requires held-out evidence?\\
Recurrence & How can the forward context exceed the truncated backward horizon?\\
LSTM & Which derivative does a product of forget gates describe?\\\bottomrule
\end{tabularx}

\subsection*{Notation to keep consistent}
$x_t$ is a token ID; a state after consuming $x_t$ predicts $x_{t+1}$.
$K_{\mathrm{in}}$ counts input rows and $K$ counts possible outputs.
$B,T,d,m,H$ denote batch size, input positions, embedding width, fixed
context length, and recurrent/hidden width. $E$ is the input table;
readout matrices are identified beside their equations. In the LSTM cell,
the subscript $o$ names its output gate rather than the vocabulary readout.
$\log$ means natural logarithm; $\log_2$ measures bits. MiB means $2^{20}$ bytes.

\subsection*{Course sources and numerical provenance}
The source baseline is repository commit \texttt{1adbaa5}. The
\href{https://github.com/baojian/llm-26-fall/tree/1adbaa5442a567cbd2ab65035b8e485778dd46a9/slides/lecture-01}{L01},
\href{https://github.com/baojian/llm-26-fall/tree/1adbaa5442a567cbd2ab65035b8e485778dd46a9/slides/lecture-02}{L02},
\href{https://github.com/baojian/llm-26-fall/tree/1adbaa5442a567cbd2ab65035b8e485778dd46a9/slides/lecture-03}{L03}, and
\href{https://github.com/baojian/llm-26-fall/tree/1adbaa5442a567cbd2ab65035b8e485778dd46a9/slides/lecture-04}{L04}
directories contain the slides, teaching plans, and classroom notebooks.
The revised notes expand recurrent models and LSTMs at the instructor's
request and place attention in Lecture 05; the older L04 deck has a different boundary.

\begin{itemize}
\item In the lecture directories, F07 redraws L01's
\path{assets/heldout-chinese.json}; F15 uses L03's \path{assets/counts-ppmi.json}.
\item F23 redraws the two distinct \texttt{tiny-lm-loss.json} assets in
L03 and L04. Their captions state models, batches, seeds, and update budgets.
\item The L02 loop numbers come from \texttt{self-training-loop.json}.
Full settings and data fingerprints are in \texttt{lecture02-results.json}
and its asset README; upstream dataset revisions were not recorded for the original download.
\item Other numerical illustrations are exact arithmetic or explicitly
chosen toy values. No new model training is reported by these notes.
\end{itemize}

The next four pages give the selected 15-paper reading route. Supporting
sources for gradient clipping and the later forget gate are linked in
Sections 7--8. The subsequent Transformer paper is
\href{https://arxiv.org/abs/1706.03762}{Vaswani et al.\ (2017)}; its
architecture belongs to the next part of the course.

\setcounter{secnumdepth}{0}
% Included by outline.tex; entries are formatted from references.bib.
\clearpage
\section{Reading guide: 15 papers before the Transformer}
These readings follow the teaching sequence, rather than publication order.
\textbf{Core} marks six recommended close readings: \readingref{R03},
\readingref{R04}, \readingref{R06}, \readingref{R10}, \readingref{R13}, and
\readingref{R14}. For the other papers, use the selected ideas below.
The labels guide study depth; they do not add graded requirements.

\textbf{Boundary.} \readingref{R01}--\readingref{R13} support the first-four
recap. \readingref{R14}--\readingref{R15} prepare Lecture 05, where attention
begins. The bibliography is supporting material outside the proposed
20--25 core pages.

\subsection{1. Statistical language modeling and tokenization}
\paperentry{R01}{Shannon (1948)}{Selected passages}
  {Sections 1, 3--4 \quad / \quad F08, F11--F12}
  {shannon1948}
  {Connect probabilistic descriptions of text, uncertainty, and compression.
  This gives the later likelihood and perplexity discussion its information-theoretic setting.}
  {The successive approximations to English and the entropy definition.
  Recompute entropy for a tiny distribution; leave communication-channel
  proofs outside this recap.}

\paperentry{R02}{Kneser and Ney (1995)}{Selected idea}
  {Section 3 \quad / \quad F09--F10}
  {kneser1995}
  {Sparse counts need a useful fallback distribution. A word's variety of
  preceding contexts provides information that its raw frequency misses.}
  {The motivation for continuation probabilities and the backoff construction.
  Use one small count example; a full smoothing implementation can remain
  supporting material.}

\paperentry{R03}{Sennrich et al.\ (2016)}{Core}
  {Section 2 \quad / \quad F05--F07}
  {sennrich2016}
  {Subword units reduce the rare-word problem while keeping a finite
  vocabulary. BPE supplies a concrete learned segmentation algorithm.}
  {Section 3.2 and the merge algorithm. Trace one merge, then distinguish
  the paper's character-based, within-word BPE from the course's byte-level
  tokenizer.}

\clearpage
\section{Reading guide: neural language models and embeddings}
\textbf{Thread.} Learn reusable token representations through prediction,
then compare local prediction objectives with a global co-occurrence objective.
Keep normalized language-model likelihood separate from embedding objectives.

\paperentry{R04}{Bengio et al.\ (2003)}{Core}
  {Sections 5--6 \quad / \quad F18, F20, F22}
  {bengio2003}
  {Learn word vectors and a conditional language model together. Shared
  features let the model generalize beyond individually observed word sequences.}
  {Figure 1 and Section 2. Annotate the lookup, fixed context window,
  hidden layer, and normalized output; connect their shapes to the course's
  token-to-loss computation.}

\paperentry{R05}{Mikolov et al.\ (2013a)}{Selected models}
  {Section 5 \quad / \quad F16, F18}
  {mikolov2013efficient}
  {CBOW and skip-gram make the prediction direction explicit: use context
  to predict a word, or use a word to predict surrounding words.}
  {The two architectures and their computational motivation. Draw both
  arrow patterns; use \readingref{R06} for the negative-sampling objective.}

\paperentry{R06}{Mikolov et al.\ (2013b)}{Core}
  {Section 5 \quad / \quad F16--F17}
  {mikolov2013distributed}
  {Negative sampling learns representations from observed and sampled
  word--context pairs; subsampling and phrase construction extend the approach.}
  {Section 2.2. Compute one positive and one negative term and identify
  the two embedding tables. This objective is not a normalized next-token
  log likelihood.}

\paperentry{R07}{Pennington et al.\ (2014)}{Selected objective}
  {Section 5 \quad / \quad F15--F18, conceptual comparison}
  {pennington2014}
  {GloVe learns vectors from weighted log co-occurrence counts, offering a
  useful comparison with the predictive objectives of word2vec.}
  {The co-occurrence-ratio motivation and weighted least-squares objective.
  Identify counts, weights, vectors, and biases; distinguish this objective
  from the PPMI/SVD example.}

\clearpage
\section{Reading guide: recurrence, gradients, and memory}
\textbf{Thread.} A recurrent state carries history beyond a fixed input
window. Training this state exposes long gradient paths, motivating
controlled access to a memory cell.

\paperentry{R08}{Elman (1990)}{Selected architecture}
  {Section 7 \quad / \quad F26--F27}
  {elman1990}
  {A hidden state copied into the next step lets a simple recurrent network
  represent temporal context using shared parameters.}
  {The network diagram and selected sequence examples. Unroll three steps
  and distinguish a changing hidden state from fixed, shared recurrent weights.}

\paperentry{R09}{Bengio et al.\ (1994)}{Selected argument}
  {Section 7 \quad / \quad F27--F28}
  {bengio1994}
  {Learning long dependencies with gradient descent can be difficult
  because sensitivity changes across repeated recurrent transformations.}
  {The long-dependency argument and its assumptions. Connect it to
  products of Jacobians and the scalar-power illustration; avoid treating
  that illustration as a measured RNN gradient.}

\paperentry{R10}{Hochreiter and Schmidhuber (1997)}{Core}
  {Section 8 \quad / \quad F29--F30}
  {hochreiter1997}
  {LSTM introduces a memory cell and gated access to help preserve useful
  learning signals over long intervals.}
  {The memory-cell mechanism and gradient motivation. The original cell
  predates the forget gate; the notes use the later formulation identified
  in the supporting references. Gates do not guarantee successful learning.}

\paperentry{R11}{Mikolov et al.\ (2010)}{Selected model and results}
  {Sections 6--7 \quad / \quad F20, F26}
  {mikolov2010}
  {Apply a recurrent hidden state to probabilistic language modeling,
  connecting the sequence architecture directly to next-word prediction.}
  {Figure 1, the model description, and a results table. Distinguish
  static evaluation from the paper's dynamic test-time updates before
  comparing perplexities.}

\clearpage
\section{Reading guide: encoder--decoder models and attention}
\textbf{Boundary.} The first two entries close the RNN/LSTM story in
Lecture 04. The last two begin Lecture 05. Their annotations motivate the
next lecture without adding attention mechanisms to the first-four recap.

\paperentry{R12}{Cho et al.\ (2014)}{Selected architecture}
  {Section 8 \quad / \quad F29, F32, comparison only}
  {cho2014}
  {An RNN encoder--decoder learns phrase representations, with reset and
  update gates providing another approach to controlling recurrent memory.}
  {Section 2 and the gated hidden unit. Compare its gates with LSTM and
  draw the encoder-to-decoder interface; keep a full GRU derivation as
  supporting reading.}

\paperentry{R13}{Sutskever et al.\ (2014)}{Core}
  {Section 8; endpoint of Lecture 04 \quad / \quad F31--F32}
  {sutskever2014}
  {An LSTM encoder maps an input sequence to a fixed-dimensional vector
  from which an LSTM decoder generates an output sequence.}
  {The model diagram and reversed-input experiment. Ask what the final
  encoder representation must preserve, and which computations still depend
  on earlier sequence steps.}

\paperentry{R14}{Bahdanau et al.\ (2015)}{Core: Lecture 05}
  {Lecture 05 handoff \quad / \quad follow the F32 bottleneck discussion}
  {bahdanau2015}
  {Learn a soft alignment so the decoder can form a context from encoder
  annotations at each output step, addressing the fixed-vector interface.}
  {Sections 2--3 and the alignment illustration. Identify scores, normalized
  weights, and the weighted context. Cite the 2015 conference paper and
  acknowledge its 2014 preprint.}

\paperentry{R15}{Luong et al.\ (2015)}{Selected methods: Lecture 05}
  {Lecture 05 handoff \quad / \quad compare with \readingref{R14}}
  {luong2015}
  {Global and local attention offer distinct ways to select source
  positions and compute the context used by a recurrent decoder.}
  {Section 3: dot, general, and concat scores; global versus local context.
  Distinguish these encoder--decoder attention mechanisms from the
  Transformer self-attention architecture introduced later.}

\end{document}
